\documentclass[11pt]{article}
\usepackage[utf8]{inputenc}
\usepackage[margin=1in]{geometry}
\usepackage{hyperref}
\usepackage{setspace}

\begin{document}

\noindent
Carson Slater \hfill web: \href{https://carsonslater.github.io}{https://carsonslater.github.io} \hfill carsonslater7@gmail.com \\
\rule{\textwidth}{0.4pt}

\vspace{2em}

\begin{center}
    \Large \textbf{Cover Letter: Data Science Analyst -- Mayo Clinic}
\end{center}

\vspace{1em}

\noindent
To the Data Science Hiring Committee at Mayo Clinic,

\vspace{1em}

One of my deepest convictions is that my vocation should enable me to engage with work that has a transformative, positive influence on the world. My journey into statistics began with a prevailing thought during the early days of the pandemic: \textit{``Wow, data is so powerful.''} Since then, my driving mission has been to use statistical methods to empower people with insights for good stewardship, especially within the realm of medicine. I am applying for the Data Science Analyst position because Mayo Clinic's commitment to putting the needs of the patient first perfectly aligns with my desire to build mathematical methods for social good. I want to deploy machine learning and analytic solutions to help clinicians uncover patterns in heterogeneous data that directly impact patient health.

My appetite for understanding and deriving models to map medical phenomena was cemented during the incredibly formative summer of 2022. I had the privilege of studying at the Duke and North Carolina State University Summer Institute for Biostatistics (SIBS). It was remarkably cool to immerse myself in the active research areas of biostatistics alongside top-tier faculty. There, I was exposed to the complexities of electronic health records (EHR), navigating missing data, Bayesian modeling, and survival analysis. The experience culminated in a thrilling hackathon where I partnered with peers to build a logistic regression model predicting myocardial infarction relapse. Exploring and modeling that messy, real-world data was the catalyst that convinced me to hang up my NCAA baseball spikes and commit my high-bandwidth work ethic fully to graduate studies in statistics.

Currently, as a Statistical Science PhD student at Baylor University, I continue to translate complex, messy data into actionable insights. Through my role as a Graduate Assistant, I have consulted for interdisciplinary research teams, including work for the American Holistic Nursing Association. This consultative experience directly translates to your need for an analyst who can work alongside informaticians, knowledge architects, and clinicians, presenting findings in easy-to-understand terms for non-technical users at the point of care. Technically, I am deeply equipped to tackle your business problems. I leverage Python and R to deploy a wide array of models—from deep learning (PyTorch/TensorFlow) and ensemble models (XGBoost) to Bayesian methods and spatiotemporal modeling.

In this age of information, the medical world is crying out for better tools to assist in clinical decision-making. I am hungry to join Mayo Clinic to create parsimonious, accurate, and efficient models that augment human capabilities in healthcare. Given my background of managing maximum credit hours, faculty-mentored research, and interdisciplinary consulting—all while having previously balanced a demanding varsity athletic schedule—I am confident in my ability to manage a varied workload of projects with multiple priorities.

Thank you for your time and consideration. I would welcome the opportunity to discuss how my technical toolkit and passion for medical stewardship make me a strong fit for your team.

\vspace{1em}
\noindent
Sincerely,

\vspace{1em}
\noindent
\textbf{Carson Slater} \\
PhD Student, Statistical Science \\
Baylor University

\end{document}
